\subsection{Auxiliary primes, variation of local conditions, and a polynomial for the length}
\label{sec:selmer-polynomial-descent}
Let $O$ be the integer ring of a finite extension of $\mathbb Q_\ell$, with uniformizer
$\pi$ and residue field $k$. Write $v(\ell)=1$, and use $\lambda_O$ for ordinary
composition length. The normalized length is $\lambda_O/v_\pi(\ell)$.
We first explain how the auxiliary primes for several lattices can be chosen by a
single descent, and then construct the polynomials governing the finite-level lengths.
Throughout this subsection the finite-level modules consist of cocycles satisfying
the specified linear relations. Passing to a quotient by coboundaries in the
middle of the argument would change the lengths.

\begin{lemma}[Varying a local condition]\label{lem:selmer-tilt}
Let $C$ have finite length over $O$, and let $A,B:C\to(O/\pi^h)^d$ be linear.
There is a set of at most $d$ elements of $k$ such that for every $a\in O$
outside those residue classes,
\[
 \lambda_O\ker(aA-B)\ \leq \lambda_O\{z\in\ker A:Bz\in A(C)\}.
\]
\end{lemma}
\begin{proof}
Put $C'=\sum_{u\in O}\ker(uA-B)$. On each summand $Bz=A(uz)$, so
$B(C')\subseteq A(C')$. Let $L$ be the inverse image of $A(C')$ in $O^d$.
It is free of rank at most $d$. Write $p:L\twoheadrightarrow A(C')$ for reduction.
By freeness lift $p$ to $s:L\to C'$ through $A$, then lift $Bs$ through $p$
to an endomorphism $U$ of $L$. Thus $(aA-B)s=p(a-U)$. Away from the roots
in $k$ of the characteristic polynomial of $U\bmod\pi$, $a-U$ is invertible.
Consequently $(aA-B)(C')$ contains $A(C')$. Subtracting the image lengths
from $\lambda_O C'$ bounds the length of the kernel on $C'$ by
$\lambda_O(C'\cap\ker A)$. The entire kernel of $aA-B$ belongs to $C'$,
and $B(C')\subseteq A(C)$. This proves the assertion.
\end{proof}

Fix a torsion level $\ell^M$, and put $e=v_\pi(\ell^M)$.
For a finite family of lattices of ranks $n_i$, write $Z_{i,h}(J)$ for the
cocycles modulo $\pi^h$ unramified outside the designated places and $J\cup\bar J$,
satisfying the original relations and the following new conditions.
At each $w\in J$ fix a tame inertia generator $\gamma_w$, a Frobenius
$\phi_w$ and a slope $s_{i,w}\in O$. Require
\[
 z(\gamma_w)=s_{i,w}z(\phi_w),\qquad
 z(\gamma_w^{-c})=s_{i,w}z(\phi_w^{-c}),\qquad
 g^{-c}=(cgc^{-1})^{-1}
\]
for the chosen complex conjugation $c$. Take the auxiliary primes and conjugates
off the old ramification, distinct from one another, and prime to $\ell$.
At a prospective new prime $w$, let $C_{i,h}$ allow ramification at $w,\bar w$
without imposing either new relation. The maps $A,B$ evaluate a cocycle on the
two indicated inertia and Frobenius elements, respectively. Here $d=2n_i$.

The equality $\ker A=Z_{i,h}(J)$ on $C_{i,h}$ follows directly. At either new
place the original lattice is unramified, so the cocycle restricts to a
homomorphism on inertia. Its image is an $\ell$-group. Wild inertia and the
prime-to-$\ell$ part of tame inertia map trivially, and the $\ell$-part of
tame inertia is procyclic. Vanishing on the chosen generator therefore
removes all new ramification. For a unit $a$, imposing slope $a^{-1}$ gives
the kernel of $aA-B$.

\begin{proposition}[Simultaneous descent]\label{prop:selmer-tilt-descent}
Suppose $|k|>2n_i e+1$ for every $i$. Suppose at some stage that an allowable
new prime for datum $i$ and a level $1\leq h\leq e$ satisfy
\[
 B(Z_{i,h}(J))\not\subseteq A(C_{i,h}).
\]
Here allowable means that the auxiliary data extend to the new prime with
arbitrary slope for $i$ and zero slope for every other datum.
One can add this prime and decrease
$D(J)=\sum_i\sum_{j=1}^e\lambda_O Z_{i,j}(J)$ by at least one.
Starting at a singleton of slope zero, this gives a set of size at most
$1+\sum_{i,j}\lambda_O Z_{i,j}(\varnothing)$ at which
$B(Z_{i,j}(J))\subseteq A(C_{i,j})$, for every allowable prime, index and level.
\end{proposition}
\begin{proof}
For the chosen $i$ take the union of the exceptional sets of
Lemma~\ref{lem:selmer-tilt} for all $e$ levels. At most $2n_i e$ residues
are excluded. Choose a nonzero remaining residue and a unit lift $a$.
At every level the new length is at most that of the old cocycles with
$Bz\in A(C_{i,j})$. This is a proper submodule at the indicated level,
so its length drops there. For other indices choose slope zero, which
simply removes the new ramification. All lengths are finite. For instance
each continuous cocycle at finite level, together with the given residual
representation, gives a homomorphism to a fixed finite affine group,
unramified off a fixed finite set enlarged by the ramification of the representation.
The finiteness theorem for number fields of bounded degree unramified
outside a fixed set gives finitely many possibilities
(J.~S.~Milne, \emph{Algebraic Number Theory}, Theorem~8.42~\cite{milneANT}).
Each step spends one unit of a nonnegative integer. At the initial prime
slope zero gives precisely the original cocycles. The bound follows, and
is independent of the order $\ell^N$ of the auxiliary characters.
\end{proof}

\paragraph{A single matrix for all coefficient specializations.}
Let $G=G_E$. Fix a lattice $O^n$ with continuous action $\rho$, a finite
ramification set $S$, linear homogeneous relations on cocycle values with
coefficients in $O$, and locally constant characters
$h_j:G\to\mathbb Z/\ell^N$ for $j\in J$. With exponents read in
$[0,\ell^N)$ put
\[
 R_X(g)=\rho(g)\prod_{j\in J}(1+X_j)^{h_j(g)}.
\]
For an integral point $x$ over any finite coefficient extension $O'$,
let $Z(x)$ be the locally constant solutions in $(O'/\ell^M)^n$ of
$z(gh)=z(g)+R_x(g)z(h)$, the given relations and vanishing on inertia
outside $S$. This definition makes sense for arbitrary $x$.
At $x_j=\zeta_j-1$ with $\zeta_j^{\ell^N}=1$ the matrices give the
desired twisted action.

There is a polynomial matrix $A(X)$ with finitely many columns such that
$Z(x)=\ker(A(x)\bmod\ell^M)$ in the coordinates described below.
To see that one set of columns suffices, let $G'$ be the intersection of
the kernels of $\rho\bmod\ell^M$ and all $h_j$. Every solution restricts
on $G'$ to a homomorphism. Let $H$ be the common kernel of all continuous
homomorphisms $G'\to\mathbb Z/\ell^M$ unramified outside the places above
$S$. There are finitely many by the same bounded-degree finiteness theorem
over the fixed field of $G'$. Replace $H$ by its normal core in $G$.
Every such $z$ vanishes on $H$: composing its restriction with additive
characters into $\mathbb Z/\ell^M$ detects each nonzero element of its
finite abelian target of exponent dividing $\ell^M$. The equation for
$z(gh)$ now shows $z$ factors as a function on $G/H$. Use $n[G:H]$
unknowns, one $n$-vector for each coset. Write all the equations in these
unknowns as rows over $O[X_j:j\in J]$, retaining all original pairs
$g,h$ when writing the cocycle equation. Conversely any solution of this
matrix gives a locally constant function satisfying every equation.
The polynomial ring is Noetherian, so a finite subset of the rows
generates their row module and suffices over every specialization.

\begin{lemma}[One polynomial sharp at a chosen point]\label{lem:selmer-polynomial}
For every integral point $x_0$ over a finite coefficient extension there is
a polynomial $P\in O[X_j:j\in J]$ such that
\[
 \frac{\lambda_{O'}Z(x)}{v_{\pi'}(\ell)}\leq v(P(x))
\]
over every finite coefficient extension $O'$ with uniformizer $\pi'$,
with equality at $x_0$.
\end{lemma}
\begin{proof}
If $A$ has $d$ columns, append the $d$ rows of $\ell^M I_d$.
At each specialization the minimum valuation of its maximal minors is
the normalized length of $Z(x)$. Indeed, put $A(x)$ in Smith normal
form over $O'$. Writing $t_j$ for its elementary-divisor valuations in
uniformizer units, with $t_j=\infty$ for zero positions, the kernel
modulo $\ell^M$ has length $\sum_{j=1}^d\min(t_j,v_{\pi'}(\ell^M))$.
The row lattice of the augmented matrix is the row lattice of $A(x)$
plus $\ell^M O'^d$, so its index has the very same length. Its index
ideal is generated by the maximal minors. Select a minor attaining the
minimum at $x_0$. Its valuation bounds the minimum at every other
point. Division by $v_{\pi'}(\ell)$ gives the assertion.
\end{proof}

For coefficients $K'/O'$ ($K'=\operatorname{Frac}(O')$) use the same
cocycle and linear relations. The $\ell^M$-torsion of this solution
module is exactly the finite-level module at a character point, via
$u\bmod\ell^M\mapsto u/\ell^M\bmod O'$. If its normalized length is
strictly below $M$, every solution is already killed by $\ell^M$.
Indeed the image of a continuous cocycle in its discrete target is finite
by compactness, so each solution is torsion. An element of order
greater than $\pi'^{v_{\pi'}(\ell^M)}$ would generate a cyclic
submodule whose $\ell^M$-torsion alone has normalized length $M$.
The polynomial therefore bounds the divisible-coefficient length
whenever $v(P(x))<M$.

\paragraph{Where the arithmetic enters.}
The descent proves its terminal inclusion without producing the
inequality $\lambda_O Z_{i,e}(J)\leq G_i+L_i$ from that inclusion.
Here $G_i$ is the supremum of lengths of equivariant maps from the
torsion lattice to its conjugate Tate dual and $L_i$ is the sum of
local excess lengths away from $\ell$. Section~\ref{sec:terminal-selmer} proves the implication using global duality and the reciprocity argument for these local conditions. The construction of the allowable prime characters
with the required values at one another's Frobenius elements is also
needed here. Section~\ref{sec:terminal-selmer} constructs the primes and proves the terminal bound under the stated hypotheses on the underlying Selmer data. Lemma~\ref{lem:selmer-polynomial} supplies the
polynomial once a point of small length has been found. A polynomial
sharp at the untwisted point need not be small there. The next
calculation explains how to produce a small point from the tied bound.

\paragraph{Releasing the local conditions along a line.}
Take $J$ and its distinguished prime $w_0$ as above, write $q=\ell^N$, and
assume the following local specifications. Every $w\in J$ has norm $1\bmod q$.
The anticyclotomic $h_j$ have values $h_j(\gamma_w)=\delta_{jw}$ and
$h_j(\phi_w)=\delta_{jw_0}$ for $w\ne w_0$ (the inertia prescription applies
also to $w_0$). The old $\rho$ is unramified at the new places.
For each nonzero slope at $w$, $\rho(\phi_w)$ and
$\rho(\phi_w^{-c})$ act trivially modulo $\ell^M$.
Fix one datum of rank $n$, and suppose that after releasing the zero slope
at $w_0$ the normalized untwisted length, with all the other ties retained,
is at most $\beta$. This last release costs at most
\[
 \sum_{g\in(\phi_{w_0},\phi_{w_0}^{-c})}v(\det(\rho(g)-1)),
\]
where a sum over a displayed pair counts both entries. Indeed the untwisted
tame relation gives $(\rho(g)-1)z(\gamma)=0$ modulo $\ell^M$ when $N\geq M$:
Frobenius raises tame inertia to the norm and that norm is $1\bmod\ell^M$.
Evaluation therefore embeds the quotient by the two restored ties into the
two kernels. Each kernel has normalized length at most its determinant
valuation by Smith normal form. If the Frobenius matrices agree with
prescribed comparison matrices modulo $\ell^M$, the same kernel bound
uses the comparison matrices.

For $0<u\leq1$ rational, $Nu\geq M$, take the line
$X_{w_0}=Y$ and $X_w=s_wY$ for $w\ne w_0$. Let $v(y)=u$ in a finite
coefficient extension. For any $\alpha\equiv1\pmod y$, the sum of
$\alpha^j$ for $0\leq j<\ell^N$ has valuation at least $Nu$.
Indeed, factor the sum into $N$ geometric sums of length $\ell$.
Each is congruent to $\ell$ modulo $y$, and $v(\ell)=1\geq u$.
It follows that all $(1+X_w)^q$ specialize to $1\bmod\ell^M$.
The point twist is a representation modulo $\ell^M$. Let $z$ be a solution
for that representation with all auxiliary ties removed. For
$w\ne w_0$ use the pair $\phi,\gamma$ at $w$ or its conjugate inverse.
In the affine representation $g\mapsto(t\mapsto R_x(g)t+z(g))$,
the inertia image has exponent dividing $q$: the linear part is scalar
$\alpha\equiv1\pmod y$ and the translation of its $q$th power is
$(1+\alpha+\cdots+\alpha^{q-1})z(g)=0$. It therefore factors through
the cyclic tame $\ell$-quotient. Frobenius commutes with this inertia
image since the norm is $1\bmod q$; so does its inverse. Comparing the
translations gives
\[
 ((1+y)\rho(\phi)-1)z(\gamma)=s_w y z(\phi).
\]
The hypotheses on the slope give
$s_w((1+y)\rho(\phi)-1)=s_w y I$ modulo $\ell^M$.
Thus the discrepancy $z(\gamma)-s_w z(\phi)$ is killed by
$(1+y)\rho(\phi)-1$. Evaluating all discrepancies embeds the quotient by
the retained ties into their respective kernels. With
\[
 H(Y)=\prod_{w\ne w_0}\,
       \prod_{g\in(\phi_w,\phi_w^{-c})}\det((1+Y)\rho(g)-1),
\]
the normalized cost of removing these ties is at most $v(H(y))$.

Apply Lemma~\ref{lem:selmer-polynomial} to the retained ties at the origin,
and restrict its polynomial to this line, obtaining $F(Y)$ with
$v(F(0))\leq\beta$. For any integral polynomial $f=\sum a_jY^j$ define
$v_u(f)=\min_j(v(a_j)+ju)$. Finitely many nonzero polynomials attain their
respective $v_u$ at the same point of valuation $u$: write $y=t\xi$ with
$v(t)=u$, divide each $f(tX)$ by a coefficient of minimal valuation, and
choose the nonzero residue of $\xi$ outside the roots of the resulting
nonzero residue polynomials. A finite unramified extension supplies this
residue. Since $H$ has unit leading coefficient and degree
$2n(|J|-1)$, at this point
\[
 v(F(y))+v(H(y))\leq\beta+2n(|J|-1)u .
\]
This bounds the untied normalized length at $x=(y,(s_wy)_{w\ne w_0})$.
Applying Lemma~\ref{lem:selmer-polynomial} at $x$ to the untied module now
gives the desired polynomial, with $v(P(x))$ bounded by the same quantity
and all coordinates of valuation at least $u$.

\paragraph{The group-ring polynomial.}
For an integral Stickelberger element $\Theta_J=\sum_b c_b[b]$, put
$P_J(X)=\sum_b c_b\prod_j(1+X_j)^{h_j(b)}$, using the nonnegative
representatives of the exponents. Here the characters descend to the ray
class group in which $b$ ranges. Substitution of roots of unity recovers exactly
the character evaluation. Suppose the one-prime norm relation multiplies this
evaluation by $(1-\alpha q_w^{-d})(1-\alpha^{-1}q_w^{-d})$ for each
$d$ in a fixed list of $r$ exponents, with $q_w$ the prime norm and
$\alpha=\vartheta(\phi_w)\psi(\phi_w)$ for a character of the lower level.
Assume the coefficients and $\vartheta$ are integral with $\vartheta$ unit
valued. Put $\beta'=v(\boldsymbol{1}(\Theta_\varnothing))+
2r v(\vartheta(g_0)-1)<\infty$ with $\vartheta(g_0)\ne1$.
Choose $M'>\beta'+1$, and require $Nu\geq M'$ and
$v(\vartheta(\phi_{w_0})-\vartheta(g_0))\geq M'$. At $w_0$ and the trivial
character the Euler factors are congruent modulo valuation $M'$ to
$1-\vartheta(g_0)^{\pm1}$ since $q_{w_0}\equiv1\pmod{\ell^N}$.
If $r>0$ these differences have valuation $<M'$; if $r=0$ there is no factor.
Hence the one-variable evaluation $F_0(Y)$ of $\Theta_{\{w_0\}}$ using
$h_{w_0}$ has $v(F_0(0))\leq\beta'$.

Set $Z=1+Y$ and $D=r(|J|-1)$. On the line $X_{w_0}=Y$, $X_w=0$ otherwise,
denote $P_J$ by $P_*(Y)$. At every $Z$ of order dividing $\ell^N$, successive
applications of the norm relation give
\[
 Z^D P_*(Y)= F_0(Y)
 \prod_{w\ne w_0}\prod_d
 (1-A_wq_w^{-d} Z)(Z-A_w^{-1}q_w^{-d}),
 \qquad A_w=\vartheta(\phi_w).
\]
Indeed $h_{w_0}(\phi_w)=1$ at these primes. Subtracting the right side
from the left yields an integral polynomial divisible by $Z^{\ell^N}-1$,
by monic division and the $\ell^N$ distinct roots in characteristic zero.
The product on the right has degree $2D$ and unit leading coefficient.
At a point $v(y)=u$ generic for this product and $F_0$, its value times
$F_0(y)$ has valuation at most $\beta'+2Du$. For $2Du<1$ this
is strictly below $M'$. Meanwhile $v((1+y)^{\ell^N}-1)\geq Nu$ by the
geometric sum estimate above. The congruence therefore gives
$v(P_*(y))\leq\beta'+2Du$, since $1+y$ is a unit.
This supplies the required witness on the same auxiliary set for each system.

Thus, given allowable prime data and the terminal Selmer length estimate,
one constant $s_0\geq2|J|\max(1,\max_i n_i,\max r)$ controls the number
of variables and all witness losses. Proposition~\ref{prop:selmer-tilt-descent}
makes its choice uniform in $N$ when those prime data exist for every $N$.
One may enlarge $N$ after fixing $u$ to satisfy both torsion thresholds.
The individual witnesses may differ. The common roots-of-unity specialization
in the preceding argument applies to the polynomials so obtained.
The construction of the Galois representations and the identification of the
actual critical-value ratios with group-ring elements satisfying the displayed
norm relation remain separate arithmetic obligations.
